% Auto-generated by build_annotation_scheme_Q4_table.py; do not edit by hand.
% Requires: booktabs, longtable, array
%
% % Auto-generated by build_annotation_scheme_Q4_table.py; do not edit by hand.
% Requires: booktabs, longtable, array
%
% % Auto-generated by build_annotation_scheme_Q4_table.py; do not edit by hand.
% Requires: booktabs, longtable, array
%
% % Auto-generated by build_annotation_scheme_Q4_table.py; do not edit by hand.
% Requires: booktabs, longtable, array
%
% \input{analysis/manuscript_tables/tex/annotation_scheme_Q4.tex}
%
% Reference examples: \ref{M1}, \ref{M9}, \ref{M13}

{\footnotesize
\setlength{\tabcolsep}{4pt}
\renewcommand{\arraystretch}{1.75}
\onecolumn
\begin{longtable}{@{}>{\raggedright\arraybackslash}p{0.06\textwidth} >{\raggedright\arraybackslash}p{0.14\textwidth} >{\raggedright\arraybackslash}p{0.30\textwidth} >{\raggedright\arraybackslash}p{0.24\textwidth} >{\raggedright\arraybackslash}p{0.24\textwidth}@{}}
        \toprule
        \textbf{Code} & \textbf{Aspect} & \textbf{Definition (AI-pointing)} & \textbf{Human-pointing examples} \\
        \midrule
\endfirsthead
        \multicolumn{5}{c}{\textbf{Table \thetable{} -- continued from previous page}} \\
        \toprule
        \textbf{Code} & \textbf{Aspect} & \textbf{Definition (AI-pointing)} & \textbf{Human-pointing examples} \\
        \midrule
\endhead
        \midrule
        \multicolumn{5}{r}{\textit{Continued on next page}} \\
        \midrule
\endfoot
        \bottomrule
        \noalign{\vskip 0.75em}
        \caption{Annotation scheme for participant comments on Compar. Q7 (justifying decision for which translation is percieved as machine translation).}
        \label{tab:annotation-scheme-q4} \\
\endlastfoot
        \codelabel{M1} & Literalism / source-bleed &
        Translation feels word-for-word or dictionary-like; source-language structure, idiom, or convention bleeds through; the method feels direct and unedited rather than adapted for English. &
        translated in a way that was corrected to make sense for English &
         \\[5pt]
        \codelabel{M2} & Word choice &
        A specific word or short phrase reads as off; wrong, stilted, archaic, jargon, bland, generic, or conversely vivid and colloquial. The reader points at the word itself. &
        in a frenzy' more accurate than 'with a fever' &
         \\[5pt]
        \codelabel{M3} & Sentence flow \& connection &
        How sentences are shaped and how they connect. Choppy, truncated, run-on, abrupt, copy-paste structure, sentences that don't link. Not grammar (use M4). &
        more variation in sentence lengths &
         \\[5pt]
        \codelabel{M4} & Grammar / tense / agreement &
        Tense inconsistency, verb agreement errors, modal misuse, overuse of grammatical forms (perfect tense, conditional ``would''). Grammar mechanics that aren't sentence shape (M3) or meaning (M10). &
        (rare --- usually only flagged when wrong) &
         \\[5pt]
        \codelabel{M5} & Formatting \& surface artifacts &
        Visible page-level mechanics not folk-theory framed: punctuation, scene dividers, quotation marks, dialogue formatting, encoding. Em-dashes here only when cited as plain observation; em-dashes as ``AI sign'' $\rightarrow$ M13. &
        slightly better formatting with proper nouns &
         \\[5pt]
        \codelabel{M6} & Wrong tone for the work &
        Whole passage is in a wrong tone for a novel, i.e., manual, textbook, technical, Wikipedia-mode; ``not how a novel would read.'' Dominant mode is inappropriate. &
        older style consistent with fantasy &
         \\[5pt]
        \codelabel{M7} & Mixed / inconsistent registers &
        Within a single passage, register switches or mixes, e.g., US/UK slang together, source-culture and target-culture mixing, formal-informal swings. Internal inconsistency, not a dominant wrong key. &
        (rare) &
         \\[5pt]
        \codelabel{M8} & Wordiness / over-explaining &
        Text uses more words than needed, i.e., padding, extra words, info-dumping, lack of economy. Different from M6; wordiness can occur in any mode. &
        economical; tight prose &
         \\[5pt]
        \codelabel{M9} & Voice \& emotional life &
        Whether prose is alive, e.g., engages, immerses, evokes mood, conveys emotion, gives characters distinct voices. Fires on flatness, mechanical feel, characters sounding alike, failure to evoke scene, inability to picture/feel. Immersion lives here. &
        made me feel like a sassy idol &
         \\[5pt]
        \codelabel{M10} & Meaning errors &
        Text gets meaning, logic, embodied physics, or factual content wrong, e.g., impossible images, lost nuance, dumbed-down content. Error is in the text, not the reader's parsing. &
        HT treats speaking underwater as hypothetical rather than given &
         \\[5pt]
        \codelabel{M11} & Creative success &
        Translation makes a non-literal creative move the reader treats as a human signature, e.g., rhyme rendering, wordplay, creative password/name, right idiom for cultural concept. AI-pointing only by absence-by-contrast. &
        password 'Make\_Me\_Famous75!' --- wouldn't think AI capable of that leap &
         \\[5pt]
        \codelabel{M12} & Comprehension difficulty &
        Reader couldn't follow, i.e., had to reread, needed the other version, confused by a specific passage. About reader experience, not text defect. &
        version of the sentence clearer in HT &
         \\[5pt]
        \codelabel{M13} & Folk theory invoked &
        Reader explicitly references what AI does, e.g., named stereotype, prior experience, capability claim. Triggers on language like ``classic sign of AI,'' ``AI tends to,'' ``I'd expect from AI,'' ``AI texts I've worked with,'' ``AI doesn't [X],'' ``beyond what AI is capable of.'' &
        Human translators manage fluidity AI doesn't match &
         \\[5pt]
\end{longtable}
\twocolumn
}

%
% Reference examples: \ref{M1}, \ref{M9}, \ref{M13}

{\footnotesize
\setlength{\tabcolsep}{4pt}
\renewcommand{\arraystretch}{1.75}
\onecolumn
\begin{longtable}{@{}>{\raggedright\arraybackslash}p{0.06\textwidth} >{\raggedright\arraybackslash}p{0.14\textwidth} >{\raggedright\arraybackslash}p{0.30\textwidth} >{\raggedright\arraybackslash}p{0.24\textwidth} >{\raggedright\arraybackslash}p{0.24\textwidth}@{}}
        \toprule
        \textbf{Code} & \textbf{Aspect} & \textbf{Definition (AI-pointing)} & \textbf{Human-pointing examples} \\
        \midrule
\endfirsthead
        \multicolumn{5}{c}{\textbf{Table \thetable{} -- continued from previous page}} \\
        \toprule
        \textbf{Code} & \textbf{Aspect} & \textbf{Definition (AI-pointing)} & \textbf{Human-pointing examples} \\
        \midrule
\endhead
        \midrule
        \multicolumn{5}{r}{\textit{Continued on next page}} \\
        \midrule
\endfoot
        \bottomrule
        \noalign{\vskip 0.75em}
        \caption{Annotation scheme for participant comments on Compar. Q7 (justifying decision for which translation is percieved as machine translation).}
        \label{tab:annotation-scheme-q4} \\
\endlastfoot
        \codelabel{M1} & Literalism / source-bleed &
        Translation feels word-for-word or dictionary-like; source-language structure, idiom, or convention bleeds through; the method feels direct and unedited rather than adapted for English. &
        translated in a way that was corrected to make sense for English &
         \\[5pt]
        \codelabel{M2} & Word choice &
        A specific word or short phrase reads as off; wrong, stilted, archaic, jargon, bland, generic, or conversely vivid and colloquial. The reader points at the word itself. &
        in a frenzy' more accurate than 'with a fever' &
         \\[5pt]
        \codelabel{M3} & Sentence flow \& connection &
        How sentences are shaped and how they connect. Choppy, truncated, run-on, abrupt, copy-paste structure, sentences that don't link. Not grammar (use M4). &
        more variation in sentence lengths &
         \\[5pt]
        \codelabel{M4} & Grammar / tense / agreement &
        Tense inconsistency, verb agreement errors, modal misuse, overuse of grammatical forms (perfect tense, conditional ``would''). Grammar mechanics that aren't sentence shape (M3) or meaning (M10). &
        (rare --- usually only flagged when wrong) &
         \\[5pt]
        \codelabel{M5} & Formatting \& surface artifacts &
        Visible page-level mechanics not folk-theory framed: punctuation, scene dividers, quotation marks, dialogue formatting, encoding. Em-dashes here only when cited as plain observation; em-dashes as ``AI sign'' $\rightarrow$ M13. &
        slightly better formatting with proper nouns &
         \\[5pt]
        \codelabel{M6} & Wrong tone for the work &
        Whole passage is in a wrong tone for a novel, i.e., manual, textbook, technical, Wikipedia-mode; ``not how a novel would read.'' Dominant mode is inappropriate. &
        older style consistent with fantasy &
         \\[5pt]
        \codelabel{M7} & Mixed / inconsistent registers &
        Within a single passage, register switches or mixes, e.g., US/UK slang together, source-culture and target-culture mixing, formal-informal swings. Internal inconsistency, not a dominant wrong key. &
        (rare) &
         \\[5pt]
        \codelabel{M8} & Wordiness / over-explaining &
        Text uses more words than needed, i.e., padding, extra words, info-dumping, lack of economy. Different from M6; wordiness can occur in any mode. &
        economical; tight prose &
         \\[5pt]
        \codelabel{M9} & Voice \& emotional life &
        Whether prose is alive, e.g., engages, immerses, evokes mood, conveys emotion, gives characters distinct voices. Fires on flatness, mechanical feel, characters sounding alike, failure to evoke scene, inability to picture/feel. Immersion lives here. &
        made me feel like a sassy idol &
         \\[5pt]
        \codelabel{M10} & Meaning errors &
        Text gets meaning, logic, embodied physics, or factual content wrong, e.g., impossible images, lost nuance, dumbed-down content. Error is in the text, not the reader's parsing. &
        HT treats speaking underwater as hypothetical rather than given &
         \\[5pt]
        \codelabel{M11} & Creative success &
        Translation makes a non-literal creative move the reader treats as a human signature, e.g., rhyme rendering, wordplay, creative password/name, right idiom for cultural concept. AI-pointing only by absence-by-contrast. &
        password 'Make\_Me\_Famous75!' --- wouldn't think AI capable of that leap &
         \\[5pt]
        \codelabel{M12} & Comprehension difficulty &
        Reader couldn't follow, i.e., had to reread, needed the other version, confused by a specific passage. About reader experience, not text defect. &
        version of the sentence clearer in HT &
         \\[5pt]
        \codelabel{M13} & Folk theory invoked &
        Reader explicitly references what AI does, e.g., named stereotype, prior experience, capability claim. Triggers on language like ``classic sign of AI,'' ``AI tends to,'' ``I'd expect from AI,'' ``AI texts I've worked with,'' ``AI doesn't [X],'' ``beyond what AI is capable of.'' &
        Human translators manage fluidity AI doesn't match &
         \\[5pt]
\end{longtable}
\twocolumn
}

%
% Reference examples: \ref{M1}, \ref{M9}, \ref{M13}

{\footnotesize
\setlength{\tabcolsep}{4pt}
\renewcommand{\arraystretch}{1.75}
\onecolumn
\begin{longtable}{@{}>{\raggedright\arraybackslash}p{0.06\textwidth} >{\raggedright\arraybackslash}p{0.14\textwidth} >{\raggedright\arraybackslash}p{0.30\textwidth} >{\raggedright\arraybackslash}p{0.24\textwidth} >{\raggedright\arraybackslash}p{0.24\textwidth}@{}}
        \toprule
        \textbf{Code} & \textbf{Aspect} & \textbf{Definition (AI-pointing)} & \textbf{Human-pointing examples} \\
        \midrule
\endfirsthead
        \multicolumn{5}{c}{\textbf{Table \thetable{} -- continued from previous page}} \\
        \toprule
        \textbf{Code} & \textbf{Aspect} & \textbf{Definition (AI-pointing)} & \textbf{Human-pointing examples} \\
        \midrule
\endhead
        \midrule
        \multicolumn{5}{r}{\textit{Continued on next page}} \\
        \midrule
\endfoot
        \bottomrule
        \noalign{\vskip 0.75em}
        \caption{Annotation scheme for participant comments on Compar. Q7 (justifying decision for which translation is percieved as machine translation).}
        \label{tab:annotation-scheme-q4} \\
\endlastfoot
        \codelabel{M1} & Literalism / source-bleed &
        Translation feels word-for-word or dictionary-like; source-language structure, idiom, or convention bleeds through; the method feels direct and unedited rather than adapted for English. &
        translated in a way that was corrected to make sense for English &
         \\[5pt]
        \codelabel{M2} & Word choice &
        A specific word or short phrase reads as off; wrong, stilted, archaic, jargon, bland, generic, or conversely vivid and colloquial. The reader points at the word itself. &
        in a frenzy' more accurate than 'with a fever' &
         \\[5pt]
        \codelabel{M3} & Sentence flow \& connection &
        How sentences are shaped and how they connect. Choppy, truncated, run-on, abrupt, copy-paste structure, sentences that don't link. Not grammar (use M4). &
        more variation in sentence lengths &
         \\[5pt]
        \codelabel{M4} & Grammar / tense / agreement &
        Tense inconsistency, verb agreement errors, modal misuse, overuse of grammatical forms (perfect tense, conditional ``would''). Grammar mechanics that aren't sentence shape (M3) or meaning (M10). &
        (rare --- usually only flagged when wrong) &
         \\[5pt]
        \codelabel{M5} & Formatting \& surface artifacts &
        Visible page-level mechanics not folk-theory framed: punctuation, scene dividers, quotation marks, dialogue formatting, encoding. Em-dashes here only when cited as plain observation; em-dashes as ``AI sign'' $\rightarrow$ M13. &
        slightly better formatting with proper nouns &
         \\[5pt]
        \codelabel{M6} & Wrong tone for the work &
        Whole passage is in a wrong tone for a novel, i.e., manual, textbook, technical, Wikipedia-mode; ``not how a novel would read.'' Dominant mode is inappropriate. &
        older style consistent with fantasy &
         \\[5pt]
        \codelabel{M7} & Mixed / inconsistent registers &
        Within a single passage, register switches or mixes, e.g., US/UK slang together, source-culture and target-culture mixing, formal-informal swings. Internal inconsistency, not a dominant wrong key. &
        (rare) &
         \\[5pt]
        \codelabel{M8} & Wordiness / over-explaining &
        Text uses more words than needed, i.e., padding, extra words, info-dumping, lack of economy. Different from M6; wordiness can occur in any mode. &
        economical; tight prose &
         \\[5pt]
        \codelabel{M9} & Voice \& emotional life &
        Whether prose is alive, e.g., engages, immerses, evokes mood, conveys emotion, gives characters distinct voices. Fires on flatness, mechanical feel, characters sounding alike, failure to evoke scene, inability to picture/feel. Immersion lives here. &
        made me feel like a sassy idol &
         \\[5pt]
        \codelabel{M10} & Meaning errors &
        Text gets meaning, logic, embodied physics, or factual content wrong, e.g., impossible images, lost nuance, dumbed-down content. Error is in the text, not the reader's parsing. &
        HT treats speaking underwater as hypothetical rather than given &
         \\[5pt]
        \codelabel{M11} & Creative success &
        Translation makes a non-literal creative move the reader treats as a human signature, e.g., rhyme rendering, wordplay, creative password/name, right idiom for cultural concept. AI-pointing only by absence-by-contrast. &
        password 'Make\_Me\_Famous75!' --- wouldn't think AI capable of that leap &
         \\[5pt]
        \codelabel{M12} & Comprehension difficulty &
        Reader couldn't follow, i.e., had to reread, needed the other version, confused by a specific passage. About reader experience, not text defect. &
        version of the sentence clearer in HT &
         \\[5pt]
        \codelabel{M13} & Folk theory invoked &
        Reader explicitly references what AI does, e.g., named stereotype, prior experience, capability claim. Triggers on language like ``classic sign of AI,'' ``AI tends to,'' ``I'd expect from AI,'' ``AI texts I've worked with,'' ``AI doesn't [X],'' ``beyond what AI is capable of.'' &
        Human translators manage fluidity AI doesn't match &
         \\[5pt]
\end{longtable}
\twocolumn
}

%
% Reference examples: \ref{M1}, \ref{M9}, \ref{M13}

{\footnotesize
\setlength{\tabcolsep}{4pt}
\renewcommand{\arraystretch}{1.75}
\onecolumn
\begin{longtable}{@{}>{\raggedright\arraybackslash}p{0.06\textwidth} >{\raggedright\arraybackslash}p{0.14\textwidth} >{\raggedright\arraybackslash}p{0.30\textwidth} >{\raggedright\arraybackslash}p{0.24\textwidth} >{\raggedright\arraybackslash}p{0.24\textwidth}@{}}
        \toprule
        \textbf{Code} & \textbf{Aspect} & \textbf{Definition (AI-pointing)} & \textbf{Human-pointing examples} \\
        \midrule
\endfirsthead
        \multicolumn{5}{c}{\textbf{Table \thetable{} -- continued from previous page}} \\
        \toprule
        \textbf{Code} & \textbf{Aspect} & \textbf{Definition (AI-pointing)} & \textbf{Human-pointing examples} \\
        \midrule
\endhead
        \midrule
        \multicolumn{5}{r}{\textit{Continued on next page}} \\
        \midrule
\endfoot
        \bottomrule
        \noalign{\vskip 0.75em}
        \caption{Annotation scheme for participant comments on Compar. Q7 (justifying decision for which translation is percieved as machine translation).}
        \label{tab:annotation-scheme-q4} \\
\endlastfoot
        \codelabel{M1} & Literalism / source-bleed &
        Translation feels word-for-word or dictionary-like; source-language structure, idiom, or convention bleeds through; the method feels direct and unedited rather than adapted for English. &
        translated in a way that was corrected to make sense for English &
         \\[5pt]
        \codelabel{M2} & Word choice &
        A specific word or short phrase reads as off; wrong, stilted, archaic, jargon, bland, generic, or conversely vivid and colloquial. The reader points at the word itself. &
        in a frenzy' more accurate than 'with a fever' &
         \\[5pt]
        \codelabel{M3} & Sentence flow \& connection &
        How sentences are shaped and how they connect. Choppy, truncated, run-on, abrupt, copy-paste structure, sentences that don't link. Not grammar (use M4). &
        more variation in sentence lengths &
         \\[5pt]
        \codelabel{M4} & Grammar / tense / agreement &
        Tense inconsistency, verb agreement errors, modal misuse, overuse of grammatical forms (perfect tense, conditional ``would''). Grammar mechanics that aren't sentence shape (M3) or meaning (M10). &
        (rare --- usually only flagged when wrong) &
         \\[5pt]
        \codelabel{M5} & Formatting \& surface artifacts &
        Visible page-level mechanics not folk-theory framed: punctuation, scene dividers, quotation marks, dialogue formatting, encoding. Em-dashes here only when cited as plain observation; em-dashes as ``AI sign'' $\rightarrow$ M13. &
        slightly better formatting with proper nouns &
         \\[5pt]
        \codelabel{M6} & Wrong tone for the work &
        Whole passage is in a wrong tone for a novel, i.e., manual, textbook, technical, Wikipedia-mode; ``not how a novel would read.'' Dominant mode is inappropriate. &
        older style consistent with fantasy &
         \\[5pt]
        \codelabel{M7} & Mixed / inconsistent registers &
        Within a single passage, register switches or mixes, e.g., US/UK slang together, source-culture and target-culture mixing, formal-informal swings. Internal inconsistency, not a dominant wrong key. &
        (rare) &
         \\[5pt]
        \codelabel{M8} & Wordiness / over-explaining &
        Text uses more words than needed, i.e., padding, extra words, info-dumping, lack of economy. Different from M6; wordiness can occur in any mode. &
        economical; tight prose &
         \\[5pt]
        \codelabel{M9} & Voice \& emotional life &
        Whether prose is alive, e.g., engages, immerses, evokes mood, conveys emotion, gives characters distinct voices. Fires on flatness, mechanical feel, characters sounding alike, failure to evoke scene, inability to picture/feel. Immersion lives here. &
        made me feel like a sassy idol &
         \\[5pt]
        \codelabel{M10} & Meaning errors &
        Text gets meaning, logic, embodied physics, or factual content wrong, e.g., impossible images, lost nuance, dumbed-down content. Error is in the text, not the reader's parsing. &
        HT treats speaking underwater as hypothetical rather than given &
         \\[5pt]
        \codelabel{M11} & Creative success &
        Translation makes a non-literal creative move the reader treats as a human signature, e.g., rhyme rendering, wordplay, creative password/name, right idiom for cultural concept. AI-pointing only by absence-by-contrast. &
        password 'Make\_Me\_Famous75!' --- wouldn't think AI capable of that leap &
         \\[5pt]
        \codelabel{M12} & Comprehension difficulty &
        Reader couldn't follow, i.e., had to reread, needed the other version, confused by a specific passage. About reader experience, not text defect. &
        version of the sentence clearer in HT &
         \\[5pt]
        \codelabel{M13} & Folk theory invoked &
        Reader explicitly references what AI does, e.g., named stereotype, prior experience, capability claim. Triggers on language like ``classic sign of AI,'' ``AI tends to,'' ``I'd expect from AI,'' ``AI texts I've worked with,'' ``AI doesn't [X],'' ``beyond what AI is capable of.'' &
        Human translators manage fluidity AI doesn't match &
         \\[5pt]
\end{longtable}
\twocolumn
}
